% !TeX root = ../../Book2.tex
\OptionalSection{高维 Poisson 求和}{b2:sec:0604}

整数格的直角坐标并非本质要求。任意全秩格都可以由可逆线性映射得到，频率侧相应出现逆转置映射。除了点的位置，还必须追踪基本区域的体积。

\subsection{格与对偶格}
\label{b2:subsec:060401}

\begin{definition}\label{b2:def:06-lattice}
若 \(A\) 是可逆实 \(d\times d\) 矩阵，则集合 \(\Lambda=A\mathbb Z^d\) 称为 \(\mathbb R^d\) 中的\term[ge]{格}[lattice]，本节均指全秩格。集合 \(A[0,1)^d\) 为一个基本区域，其体积
\[
 v(\Lambda)\coloneqq|\det A|
\]
称为格的\term[jiben-quyu-tiji]{基本区域体积}[covolume]。定义\term[duiou-ge]{对偶格}[dual lattice]为
\[
 \Lambda^*\coloneqq\{\,\xi\in\mathbb R^d:\xi\cdot v\in\mathbb Z\quad(\forall v\in\Lambda)\,\}.
\]
\end{definition}
\symidx{\(\Lambda^*\)：对偶格}
\symidx{\(v(\Lambda)\)：格的基本区域体积}

\begin{proposition}\label{b2:prop:06-dual-lattice}
有 \(\Lambda^*=A^{-\mathsf T}\mathbb Z^d\)、\((\Lambda^*)^*=\Lambda\) 和 \(v(\Lambda^*)=v(\Lambda)^{-1}\)。基本区域体积不依赖格基的选择。
\end{proposition}
\begin{proof}
条件 \(\xi\cdot An\in\mathbb Z\) 对全部整数向量 \(n\) 成立，等价于 \(A^{\mathsf T}\xi\) 的每个坐标都是整数，所以得到对偶格公式。再次取逆转置和行列式便得其余两个等式。若另一基矩阵为 \(B\)，则 \(U=A^{-1}B\) 及 \(U^{-1}\) 都是整数矩阵，因此其行列式为互为倒数的整数，只能为 \(\pm1\)。故 \(|\det B|=|\det A|\)。
\end{proof}

对角矩阵给出矩形格，边长分别为 \(a_j\) 时，对偶边长为 \(a_j^{-1}\)。非正交基也遵循逆转置，而不是把每个基向量分别除以其长度；内积的整数条件才是定义。

\Cref{b2:fig:06-dual-lattice}用一个非正交格显示逆转置的作用。箭头端点的内积满足 \(b_i\cdot b_j^*=\delta_{ij}\)，不能用逐个基向量取倒长度代替这一条件。
% 本书原创：同一欧氏比例绘出非正交格与逆转置对偶格。
\begin{figure}[htbp]
\centering
\begin{tikzpicture}[x=1cm,y=1cm,>=stealth,every node/.style={font=\small},line width=.65pt]
\begin{scope}[xshift=0.0cm]
\fill[AnalysisBlue!10] (0.00000,0.00000)--(1.50000,0.00000)--(2.00000,1.00000)--(0.50000,1.00000)--cycle;
\draw[->,gray] (-1.2,0)--(3.4,0);\draw[->,gray] (0,-1.2)--(0,2.6);
\fill[gray] (-1.00000,1.00000) circle(1.15pt);
\fill[gray] (-0.50000,2.00000) circle(1.15pt);
\fill[gray] (-0.50000,-1.00000) circle(1.15pt);
\fill[gray] (0.00000,0.00000) circle(1.15pt);
\fill[gray] (0.50000,1.00000) circle(1.15pt);
\fill[gray] (1.00000,2.00000) circle(1.15pt);
\fill[gray] (1.00000,-1.00000) circle(1.15pt);
\fill[gray] (1.50000,0.00000) circle(1.15pt);
\fill[gray] (2.00000,1.00000) circle(1.15pt);
\fill[gray] (2.50000,2.00000) circle(1.15pt);
\fill[gray] (2.50000,-1.00000) circle(1.15pt);
\fill[gray] (3.00000,0.00000) circle(1.15pt);
\draw[->,AnalysisBlue,thick] (0,0)--(1.50000,0.00000) node[above right] {\(b_1\)};
\draw[->,AnalysisBlue,thick] (0,0)--(0.50000,1.00000) node[above right] {\(b_2\)};
\node at (1,3.0) {\(\Lambda,\quad v(\Lambda)=3/2\)};\end{scope}
\begin{scope}[xshift=6.6cm]
\fill[AnalysisBlue!10] (0.00000,0.00000)--(0.66667,-0.33333)--(0.66667,0.66667)--(0.00000,1.00000)--cycle;
\draw[->,gray] (-1.2,0)--(3.4,0);\draw[->,gray] (0,-1.2)--(0,2.6);
\fill[gray] (-0.66667,-0.66667) circle(1.15pt);
\fill[gray] (-0.66667,0.33333) circle(1.15pt);
\fill[gray] (-0.66667,1.33333) circle(1.15pt);
\fill[gray] (0.00000,-1.00000) circle(1.15pt);
\fill[gray] (0.00000,0.00000) circle(1.15pt);
\fill[gray] (0.00000,1.00000) circle(1.15pt);
\fill[gray] (0.00000,2.00000) circle(1.15pt);
\fill[gray] (0.66667,-0.33333) circle(1.15pt);
\fill[gray] (0.66667,0.66667) circle(1.15pt);
\fill[gray] (0.66667,1.66667) circle(1.15pt);
\fill[gray] (1.33333,-0.66667) circle(1.15pt);
\fill[gray] (1.33333,0.33333) circle(1.15pt);
\fill[gray] (1.33333,1.33333) circle(1.15pt);
\fill[gray] (2.00000,-1.00000) circle(1.15pt);
\fill[gray] (2.00000,0.00000) circle(1.15pt);
\fill[gray] (2.00000,1.00000) circle(1.15pt);
\fill[gray] (2.00000,2.00000) circle(1.15pt);
\fill[gray] (2.66667,-0.33333) circle(1.15pt);
\fill[gray] (2.66667,0.66667) circle(1.15pt);
\fill[gray] (2.66667,1.66667) circle(1.15pt);
\draw[->,AnalysisBlue,thick] (0,0)--(0.66667,-0.33333) node[above right] {\(b_1^*\)};
\draw[->,AnalysisBlue,thick] (0,0)--(0.00000,1.00000) node[above right] {\(b_2^*\)};
\node at (1,3.0) {\(\Lambda^*,\quad v(\Lambda^*)=2/3\)};\end{scope}
\end{tikzpicture}
\caption[非正交格与对偶格]{非正交格与对偶格。左侧基矩阵为 \(\left(\begin{smallmatrix}3/2&1/2\\0&1\end{smallmatrix}\right)\)，右侧取逆转置；淡蓝色区域为相应基本平行四边形。两图采用同一长度单位，区域面积互为倒数。}
\label{b2:fig:06-dual-lattice}
\end{figure}


\subsection{欧氏空间中的求和公式}
\label{b2:subsec:060402}

\begin{theorem}[全秩格上的 Poisson 求和]\label{b2:thm:06-lattice-poisson}
若 \(f\in\mathcal S(\mathbb R^d)\)、\(\Lambda\) 为全秩格，则
\[
 \sum_{v\in\Lambda}f(v)
 =\dfrac1{v(\Lambda)}\sum_{\xi\in\Lambda^*}\widehat f(\xi).
\]
两侧都绝对收敛。
\end{theorem}
\begin{proof}
先取 \(\Lambda=\mathbb Z^d\)。逐坐标应用一维 Poisson 公式；固定其他坐标时，一维切片为 Schwartz 函数，部分 Fourier 变换在其余坐标及变换后的坐标上仍满足快速下降估计。后一事实由\chapref{b2:ch:05}的不变性证明逐个坐标作同样的积分估计得到。因此，每一步产生的混合求和与积分都受 \(C\prod_j(1+|x_j|)^{-2}\) 或对应整数级数控制，可以交换，依次得到 \(\sum_{n\in\mathbb Z^d}f(n)=\sum_{k\in\mathbb Z^d}\widehat f(k)\)。

一般情形令 \(g(x)=f(Ax)\)，则 \(g\in\mathcal S\)，且 \(\widehat g(k)=|\det A|^{-1}\cdot\widehat f(A^{-\mathsf T}k)\)。对 \(g\) 使用整数格公式，即得结论。绝对收敛由可逆映射的长度比较及任意高阶的快速下降保证。
\end{proof}

基本区域体积因子不能省略。例如在一维取 \(\Lambda=h\mathbb Z\)，对偶格为 \(h^{-1}\mathbb Z\)，公式恰好恢复 \(\sum_nf(hn)=h^{-1}\sum_k\widehat f(k/h)\)。高维中的行列式就是这个因子的体积版本。

\subsection{格的 \texorpdfstring{\(\theta\)}{\AnalysisPdfTheta} 级数}
\label{b2:subsec:060403}

\begin{definition}\label{b2:def:06-lattice-theta}
对全秩格 \(\Lambda\) 与 \(t>0\)，定义它的\term[ge-theta-jishu]{格 \(\theta\) 级数}[lattice theta series]为
\[
 \Theta_\Lambda(t)\coloneqq\sum_{v\in\Lambda}e^{-\pi t|v|^2}.
\]
\end{definition}
\symidx{\(\Theta_\Lambda\)：全秩格的 \(\theta\) 级数}

\begin{corollary}\label{b2:cor:06-lattice-theta}
对所有 \(t>0\)，
\[
 \Theta_\Lambda(t)=v(\Lambda)^{-1}t^{-d/2}\Theta_{\Lambda^*}(1/t).
\]
\end{corollary}
\begin{proof}
在格上的 Poisson 公式中取 \(f(x)=e^{-\pi t|x|^2}\)。其变换为 \(t^{-d/2}e^{-\pi|\xi|^2/t}\)，代入即得。
\end{proof}

若 \(\Lambda=\mathbb Z^d\)，则 \(\Theta_\Lambda=\vartheta^d\)。一般格要同时把参数取倒数、把格换成对偶格，并保留基本区域体积；只对自对偶格才可将两边写成同一个格的级数。若用二次型矩阵表示格，\(|An|^2=n^{\mathsf T}A^{\mathsf T}An\)，相应的频率二次型由逆矩阵给出。这为第三卷的二次型 \(\theta\) 级数提供几何来源，但不要求在这里提前引入整数性或模形式的附加结构。
